% ---------------------------------------------------------------------------
% Answers the proposal's question in the terms the proposal asked it, then says what
% supersedes the question. The order is deliberate: the proposal's answer first,
% because the proposal asked it, and the finding second, because it is larger.
% ---------------------------------------------------------------------------

\section{Conclusion}
\label{s2:sec:conclusion}

The question this work was proposed to answer was whether observability lock-in is a
syntactic accident or a semantic necessity. It has an answer, and then it has a successor.

The answer is that it depends on the direction. Leaving a proprietary vendor for an open
one is a syntactic accident: we built the translator the proposal described --- two inferred
grammars, both admitting deterministic one-token-lookahead parsers, a syntax-directed
lowering into an intermediate representation, an emission into PromQL, soundness stated
relative to an ingestion mapping because the two systems do not store the same thing ---
and it translates $96.6\%$ of the $8{,}770$ Datadog metric query occurrences in our corpus,
with a ceiling of $99.9\%$ and a shortfall that is five named constructs rather than a
residue. The column we had reserved for things PromQL cannot express emptied of everything
we can demonstrate, $14$ occurrences aside, and most of what emptied it were corrections to
our own claims. Coming back is a semantic necessity and the necessity is provable: of the
$84$ refusals in that direction, $74$ rest on a theorem about single globs together with a
branch bound we state and $10$ on a theorem about globs outright, and no implementation
effort removes those ten. Leaving costs work; returning costs expressiveness.

The successor question is the one this paper is named for, and it is larger because it bears
on the direction people actually travel. A dashboard query does not determine its own
meaning. $97.0\%$ of our translations depend on a rollup interval the source query leaves
to the dashboard's time range; where a widget pins that range the obvious default is wrong
two times in three, and for the $7{,}721$ occurrences where nothing pins it there is no
value anywhere to recover. All $8{,}475$ translations assume a metric type no query in
either language states, and against the vendor's own declarations that assumption is wrong
for $491$ panels a migration would carry across believing them exact. $88.6\%$ of Datadog
and $82.5\%$ of PromQL occurrences carry a template variable, and $87.9\%$ of PromQL tag
matchers have a variable where their pattern should be --- so for most of the corpus the
reverse direction is not hard but statically undefined. A competitor's translator, which
never saw this corpus, supplies an interval to $21.3\%$ of occurrences and says so in
prose.

What a migration ports is therefore a query schema plus a resolution layer --- time range,
variable bindings, metric metadata --- that belongs to the dashboard tool and appears in
neither language's grammar. The switching cost we set out to measure is mostly there, and
it is invisible to every instrument built to measure it: a query that does not fix its own
meaning parses like any other, translates like any other, and is counted as covered like
any other. That is why it took the whole apparatus to find, and why a reader who takes
$96.6\%$ as the practical difficulty of leaving will have taken the wrong number.

Two things follow and are not pursued here. The measurement that matters next is not of
query languages at all but of the resolution layer above them, which belongs to neither
vendor and which a common query language would not reach. And the experiment that would
settle this paper's own conclusion is the one it cannot perform: a differential run on
private dashboards in which the metric type and the rollup interval are deliberately wrong,
measuring what a migration loses when it supplies the context badly rather than not at
all.
